% !TEX root = ../../vzor.tex
%% Ukázka: jednotky SI (siunitx)
%% Vkládá ji soubor vzor/kapitoly/prakticka.tex; ve vzor.pdf je pod nadpisem uvedeno „Zdroj ukázky".

\section{Sazba jednotek SI pomocí balíčku \texttt{siunitx}}
\zdrojukazky{vzor/ukazky/jednotky.tex}

K správné sazbě SI jednotek je vhodné použít package \texttt{siunitx}.
Jednotky se pak zapisují pomocí příkazu \verb|\SI{číslo}{jednotka}|  

Kódem: 
{\footnotesize
\begin{verbatim}
\SI{1}{\metre\squared}; 
\SI[{per-mode = symbol-or-fraction}]{1}{\meter\per\second} = \SI[per-mode = symbol]{3.6}{\kilo\meter\per\hour};
\SI[{per-mode = symbol}]{2.3}{\meter\per\second\squared} 
$$
 \SI{1}{\micro\meter}\times\SI{1}{\micro\meter} = \SI{1}{\micro\meter\squared}
$$ 
$$
 \SI{1}{\milli\meter}\times\SI{1}{\micro\meter} = \SI{1e-9}{\meter\squared}\\
$$
\end{verbatim}}
vysázíme následující text

\SI{1}{\metre\squared}; 
  \SI[{per-mode = symbol-or-fraction}]{1}{\meter\per\second} = \SI[per-mode = symbol]{3.6}{\kilo\meter\per\hour};
  \SI[{per-mode = symbol}]{2.3}{\meter\per\second\squared} 
$$
   \SI{1}{\micro\meter}\times\SI{1}{\micro\meter} = \SI{1}{\micro\meter\squared}
$$ 
$$
   \SI{1}{\milli\meter}\times\SI{1}{\micro\meter} = \SI{1e-9}{\meter\squared}\\
$$

\medskip
\noindent\textit{Poznámka:} od verze~3 (2021) doporučuje balíček \texttt{siunitx} zápis
\verb|\qty{9.81}{\meter\per\second\squared}| místo \verb|\SI{...}{...}| a \verb|\unit{...}| místo \verb|\si{...}|.
Starší příkazy \verb|\SI| a \verb|\si| fungují dál, proto je tento vzor zatím používá.
